\documentclass[12pt,a4paper]{article}
\usepackage[utf8]{inputenc}
\usepackage[T2A]{fontenc}
\usepackage[russian]{babel}
\usepackage{amsmath,amsfonts,amssymb}
\usepackage{graphicx}
\usepackage[
    colorlinks=true,
    linkcolor=blue,
    citecolor=blue,
    urlcolor=blue,
    bookmarks=true,
    bookmarksnumbered=true,
    pdfstartview=FitH
]{hyperref}
\usepackage[left=2.5cm,right=2.5cm,top=2cm,bottom=2cm]{geometry}
\usepackage{indentfirst}
\usepackage{setspace}
\onehalfspacing
\usepackage{titlesec}
\titleformat{\section}{\Large\bfseries}{\thesection}{1em}{}
\titleformat{\subsection}{\large\bfseries}{\thesubsection}{1em}{}
\usepackage{enumitem}
\usepackage{float}

\title{Базовая реализация фреймворка для бенчмаркинга\\
производительности и численного качества методов PySATL Core}
\author{Мызников Фёдор Денисович}
\date{Санкт-Петербург\\2026}

\begin{document}

\begin{titlepage}
\centering
{\small Санкт-Петербургский государственный университет \\
Кафедра системного программирования}\\[1.5cm]

{\large \textbf{Базовая реализация фреймворка для бенчмаркинга \\
производительности и численного качества \\
методов PySATL Core}}\\[0.8cm]

{\normalsize Отчёт о прохождении учебной практики}\\[1.5cm]

\begin{flushleft}
Студент: \hfill Мызников Ф.Д. \\
Научный руководитель: \hfill доц., к.ф.-м.н. Гориховский В.И. \\
Консультант: \hfill Ёлкин Л.Н.
\end{flushleft}

\vfill
{\normalsize Санкт-Петербург\\2026}
\end{titlepage}
\newpage

\phantomsection
\pdfbookmark[1]{Содержание}{toc}
\tableofcontents
\newpage

\section{Введение}
\label{sec:intro}
PySATL Core содержит несколько реализаций для одних и тех же
характеристик распределения, преобразования или генерации выборки.
Модульные тесты подтверждают корректность контрактов, но не дают ответа
о практическом качестве: такие реализации могут различаться по времени
выполнения, потреблению памяти, численной точности и устойчивости
вблизи хвостов, границ и особенностей плотности. Отсутствие единой
методологии сравнения затрудняет обоснованный выбор метода и
воспроизводимое сопоставление с внешними библиотеками.

Внешние библиотеки вероятностных и численных методов (SciPy,
NumPy random, UNU.RAN) предоставляют эквивалентную функциональность,
однако их сопоставление с PySATL Core требует единого подхода к
постановке сценариев, измерению характеристик и хранению результатов.
Прямое сравнение «в лоб» некорректно из-за различий в параметризациях,
соглашениях о носителе, семантике массивов и граничном поведении.

Настоящая работа выполняется в рамках issue
«Benchmarking framework for performance and numerical quality»
и посвящена разработке базовой версии фреймворка, который делает
такое сравнение воспроизводимым и интерпретируемым.
\newpage

\section{Цель и задачи}
\label{sec:goals}

\subsection*{Формальная цель}

Разработать и реализовать базовую версию воспроизводимого фреймворка
для сравнительного анализа производительности и численного качества
нескольких реализаций одних и тех же характеристик в PySATL Core,
а также для их сопоставления с внешними вероятностными и численными
библиотеками при эквивалентности функциональности и совместимости
лицензий.

\subsection*{Задачи работы}

Обзор литературы ведётся параллельно и не выделяется в отдельную
задачу. Работа рассчитана на учебный год.

\paragraph{Осень 2026.}
\begin{enumerate}[label=\arabic*., leftmargin=*]

  \item \textbf{Спроектировать контракты и архитектуру фреймворка.}
  \begin{enumerate}[label=\arabic{enumi}.\arabic*., leftmargin=*]
    \item Формат сценария бенчмарка.
    \item Контракт адаптера библиотеки.
    \item Контракт поставщика эталона.
    \item Контракт сервиса измерений, включая раздельные
          prepare/fit и evaluate и подключаемое измерение памяти.
    \item Модель результата, набор статусов и метаданные окружения.
  \end{enumerate}

  \item \textbf{Реализовать минимальный вертикальный срез.}
  \begin{enumerate}[label=\arabic{enumi}.\arabic*., leftmargin=*]
    \item Один сценарий: Normal, операция \texttt{pdf}.
    \item Минимальный раннер: загрузка, валидация, один прогон,
          структурированное хранение результата.
    \item Генерация отчёта по одному сценарию.
  \end{enumerate}

  \item \textbf{Реализовать поставщиков эталонов.}
  \begin{enumerate}[label=\arabic{enumi}.\arabic*., leftmargin=*]
    \item Аналитические формулы для Normal, Uniform, Exponential.
    \item Статус валидности эталона.
  \end{enumerate}

  \item \textbf{Реализовать базовые метрики численного качества.}
  \begin{enumerate}[label=\arabic{enumi}.\arabic*., leftmargin=*]
    \item Абсолютная ошибка.
    \item Логарифмическая относительная ошибка (LRE) с явной политикой
          для нулевого эталона, антипереполнения и бесконечностей.
  \end{enumerate}

  \item \textbf{Реализовать smoke-режим фреймворка.}

  \item \textbf{Расширить покрытие операций и распределений.}
  \begin{enumerate}[label=\arabic{enumi}.\arabic*., leftmargin=*]
    \item Операции \texttt{cdf}, функция выживания, \texttt{ppf},
          моменты и \texttt{cf}.
    \item Распределения ContinuousUniform и Exponential.
    \item Распределение Лапласа (при наличии реализации).
    \item Квадрат Uniform(0, 1) (при наличии реализации).
  \end{enumerate}

  \item \textbf{Реализовать продвинутые метрики.}
  \begin{enumerate}[label=\arabic{enumi}.\arabic*., leftmargin=*]
    \item Равномерная норма \(\lVert\cdot\rVert_\infty\) и поточечные
          ошибки.
    \item Квантильные метрики.
    \item Развёртки сходимости по параметру численного метода.
  \end{enumerate}

  \item \textbf{Реализовать адаптер SciPy.}
  \begin{enumerate}[label=\arabic{enumi}.\arabic*., leftmargin=*]
    \item Сопоставление параметров, носителей и операций.
    \item Первые сопоставимые сценарии Core-vs-SciPy.
  \end{enumerate}

\end{enumerate}

\paragraph{Весна 2027.}
\begin{enumerate}[label=\arabic*., leftmargin=*, start=9]

  \item \textbf{Расширить набор сценариев.}
  \begin{enumerate}[label=\arabic{enumi}.\arabic*., leftmargin=*]
    \item Дискретные распределения (при наличии).
    \item Сценарии с разрывами, особенностями и граничными случаями.
  \end{enumerate}

  \item \textbf{Реализовать сценарии сэмплирования через
        преобразование интегральной вероятности.}

  \item \textbf{Реализовать измерение памяти.}

  \item \textbf{Обеспечить CI.}

  \item \textbf{Подготовить отчёт и презентацию к защите.}

\end{enumerate}

\subsection*{Критерии успеха}

\begin{itemize}
  \item Фреймворк воспроизводим: одинаковый запуск при фиксированном
        окружении даёт одинаковые выходные данные.
  \item Сценарии Core-vs-Core выполняются через публичные API
        PySATL Core и доходят до стадии отчёта.
  \item Реализованные метрики численного качества корректно работают
        в крайних случаях: при значениях, близких к нулю, на границах
        носителя и в хвостах распределения, где относительные метрики
        могут терять устойчивость.
  \item Несопоставимые, неподдерживаемые и сбойные случаи получают
        явный статус и не смешиваются с успешно измеренными.
  \item Результаты и отчёты сохраняются в структурированном виде,
        пригодном для последующего анализа и сравнения между запусками.
\end{itemize}

Расхождение измеренного значения с эталоном является результатом
бенчмаркинга, а не дефектом фреймворка. Дефектом считается ситуация,
когда сценарий не доходит до стадии отчёта из-за ошибки в самом
фреймворке, а не из-за статуса тестируемой реализации.
\newpage

\section{Ожидаемые результаты}
\label{sec:expected}
По итогам работы предполагается получить:

\begin{itemize}
  \item спецификацию контрактов сценария, адаптера, эталона, измерения
        и результата;
  \item раннер, выполняющий сценарии Core-vs-Core через публичные API
        PySATL Core и сохраняющий структурированные результаты
        и отчёты;
  \item набор начальных сценариев, покрывающих функциональные
        и численные характеристики;
  \item метрики численного качества: LRE, равномерную норму,
        поточечные ошибки и квантильные метрики;
  \item адаптер SciPy с документированными сопоставлениями и первые
        сценарии Core-vs-SciPy;
  \item сценарии непрерывного сэмплирования с проверкой через
        преобразование интегральной вероятности;
  \item раздельные измерения времени подготовки и оценки, а также
        контракт подключаемого измерения памяти;
  \item ручной CI workflow с артефактами.
\end{itemize}
\newpage

\renewcommand{\refname}{Список литературы}
\phantomsection
\addcontentsline{toc}{section}{Список литературы}
\label{sec:bibliography}
\begin{thebibliography}{99}

\bibitem{issue} Issue «Benchmarking framework for performance and
numerical quality». ---
URL: \url{https://github.com/PySATL/pysatl-core/issues/124}

\bibitem{core_repo} PySATL Core. ---
URL: \url{https://github.com/PySATL/pysatl-core}

\bibitem{hormann2004} Hörmann W., Leydold J., Derflinger G.
Automatic Nonuniform Random Variate Generation. ---
Springer, Statistics and Computing, 2004.

\bibitem{devroye1986} Devroye L.
Non-Uniform Random Variate Generation. ---
New York: Springer, 1986.

\bibitem{nist_strd} NIST Statistical Reference Datasets. ---
URL: \url{https://www.itl.nist.gov/div898/strd/}

\bibitem{scipy_stats} SciPy documentation: \texttt{scipy.stats}. ---
URL: \url{https://docs.scipy.org/doc/scipy/reference/stats.html}

\bibitem{unuran} Leydold J., Hörmann W., Janka E., Tirler G.
UNU.RAN User Manual. ---
URL: \url{https://statmath.wu.ac.at/unuran/}

\bibitem{mpmath} Johansson F. et al.
mpmath: a Python library for arbitrary-precision floating-point
arithmetic. --- URL: \url{https://mpmath.org/}

\bibitem{gha_dispatch} GitHub Actions: \texttt{workflow\_dispatch}
event. ---
URL: \url{https://docs.github.com/en/actions}

\end{thebibliography}

\end{document}
